\documentclass[11pt]{article}
\usepackage[margin=1in]{geometry}
\usepackage{amsmath,amssymb,bm,booktabs}
\usepackage[hidelinks]{hyperref}
\newcommand{\E}{\mathbb E}
\newcommand{\Var}{\operatorname{Var}}
\newcommand{\Cov}{\operatorname{Cov}}
\newcommand{\Tr}{\operatorname{Tr}}
\newcommand{\bstar}{{\beta^{\star}}}
\newcommand{\bhat}{\hat{\beta}}
\newcommand{\x}{\mathbf{x}}
\title{Accentuation from Nonzero and Multiple Seeds\\Uncentered Paths, Calibration Metrics, and the Required RMT Quantities}
\author{Theory extension for AccentuationPredRMT}
\date{\today}
\begin{document}
\maketitle
\begin{abstract}
The centered, zero-seed theory reduces own-gradient accentuation to the gain
ratio $g=(\bstar)^\top\bhat/\|\bhat\|^2$. Nonzero seeds introduce baseline errors;
uncentered steps introduce mean displacement. These are distinct effects.
We derive exact conditional-on-fit MSE, coefficient of determination,
true-on-predicted regression slope, and Pearson correlation, for one seed,
multiple seeds, and seed-dependent target schedules. For independent natural
seeds, pooled metrics combine ordinary prediction moments and the same
control gain ratio. The additional RMT objects are seed-weighted bilinear
and quadratic resolvent forms; finite-fit averages additionally require their
joint fluctuations with the gain ratio.
\end{abstract}

\section{Scope, sources, and conventions}
\begingroup\sloppy
Notation follows the main manuscript
\path{accentuation_rmt_main_iclr2026.tex} and its \path{math_commands.tex}
in \href{https://github.com/Animadversio/AccentuationPredictionRMTDissociation}{AccentuationPredictionRMTDissociation}.
This note extends the centered-$\alpha$ derivation in
\path{~/Writings/AccentuationPredictionRMTDissociation/linear_reg_weight_deviation.tex},
and the repository notes
\path{notes/dissociation_prediction_control_feature_space.tex} and
\path{notes/accentuation_varR_theory.tex}.
The identities below are direct algebraic derivations; the RMT substitutions
are explicitly separated from them. Training remains the centered linear
ridge problem of those sources. A nonzero evaluation seed does not require
noncentered training data.
\par\endgroup

Condition on one fitted pixel weight $\bhat$ (or
$\bhat=F^\top\hat w$ for fixed linear features). Allow constant
intercepts for completeness:
\[
 P(\x)=\widehat c+\bhat^\top\x,\qquad
 Y(\x)=c_\star+(\bstar)^\top\x,\qquad
 \bm e=\bhat-\bstar,\quad \delta c=\widehat c-c_\star.
\]
The existing centered theory sets both intercepts to zero. Define
\[
 \x=\x_0+\alpha\bhat,\qquad
 D=\bhat^\top\bhat>0,\quad N=(\bstar)^\top\bhat,\quad g=N/D,\quad L=D-N.
\]
Here $\alpha$ is the accentuation step, \emph{not} the software ridge penalty;
the latter is denoted $\alpha_{\mathrm{Ridge}}=n\lambda$. All evaluation expectations condition on
the fit unless explicitly subscripted by ``fit.'' Teacher responses are
noiseless. Slope always means regression of $Y$ on $P$ with a fitted intercept.
Our metrics are
\begin{equation}\label{eq:metrics}
 E=\E(P-Y)^2,\quad R^2=1-\frac{E}{\Var Y},\quad
 \mathrm{slope}_p=\frac{\Cov(P,Y)}{\Var P},\quad
 r_p=\frac{\Cov(P,Y)}{\sqrt{\Var P\Var Y}}.
\end{equation}
Thus $R^2$ evaluates the original predictions, without post-hoc calibration;
it is not $r_p^2$. Undefined zero-variance cases are not assigned a value.


\paragraph{Notation aligned with the main manuscript.}
We retain $\bstar,\bhat,\x,\Sigma,\hat\Sigma,n,d,\lambda,\kappa,N,D$
and reserve $F,\Sigma_F,\hat w,w^\star$ for linear features.
Here $P(\x)=f_{\bhat}(\x)$ and $Y(\x)=f_\star(\x)$ abbreviate the
student and noiseless teacher responses; $p_0,y_0$ are their seed values.
Intercepts are an optional extension and are zero in the main manuscript.
We use $\mu_p,\Sigma_p$ for the evaluation input mean and covariance,
and $\mu_0,\Sigma_0$ for seed moments. The generic loss $E$ is
$E_{acc}$ on an accentuation distribution and $E_{gen}$ on natural inputs;
likewise $R^2$ is $R^2_{acc}$ or $R^2_{gen}$. The subscript $p$ in
$\mathrm{slope}_p,r_p$ identifies the evaluation distribution, with $r_p$
reserved for Pearson correlation. Only for compact extension formulas we set
\[
 g:=N/D=\mathrm{slope}_{acc}^{\rm single\ path},\qquad
 z_{acc}=D/N-1=g^{-1}-1\quad(N\ne0).
\]
Thus $g$ is a path gain, not Pearson $r_p$ or a pooled regression slope.
The new seed-residual symbol $A_{\rm seed}$ is a scalar, not a feature map.
All traces in this note are ordinary unnormalized $\Tr$; the main
manuscript's normalized trace is $\mathrm{tr}(T)=\Tr(T)/\dim(T)$.

\section{One nonzero seed and an uncentered step}
Let $\x_0$ be fixed and set
\[
 p_0=P(\x_0),\quad y_0=Y(\x_0),\quad
 \delta_0=p_0-y_0,\quad \E\alpha=m,\quad \Var\alpha=v>0.
\]
Since $P=p_0+D\alpha$ and $Y=y_0+N\alpha$,
\begin{equation}\boxed{
 E=vL^2+(\delta_0+mL)^2,\qquad
 R^2=1-\frac{vL^2+(\delta_0+mL)^2}{vN^2}.}
 \label{eq:single}
\end{equation}
For $N\ne0$,
\begin{equation}\boxed{
 \mathrm{slope}_p=g,\qquad a=y_0-g p_0,\qquad r_p=\operatorname{sign}(N),}
\end{equation}
where $a$ is the regression intercept. The means of the step and the seed
change calibration error but not this slope or correlation. Pearson
correlation is therefore blind to gain magnitude and baseline bias on a
single noiseless linear path. For example, $g=0.1$ gives $r_p=1$ but,
even with zero baseline and centered steps, $R^2=-80$.
A through-origin fit instead gives
\[
 \mathrm{slope}_{\rm origin}=
 \frac{vDN+(p_0+mD)(y_0+mN)}{vD^2+(p_0+mD)^2};
\]
it should not be confused with the gain ratio.

\paragraph{Nonzero seed is not equivalent to uncentered $\alpha$.}
Writing $\alpha=m+\widetilde\alpha$ replaces the seed by
$\x_0+m\bhat$, so a mean shift can be absorbed into the component of the
seed along $\bhat$. This is an exact reparameterization of the same generated images: every
uncentered step can be absorbed into a shifted seed. A general seed, however,
has a perpendicular component that cannot be absorbed into the scalar step. Conversely, a nonzero seed need not change
$\E\alpha$. The expression $\delta_0+mL$ displays both effects explicitly.


\paragraph{Does the shifted seed remain natural?}
No natural-distribution assumption is needed for the exact path identities.
It is an additional experimental assumption, useful when accentuation starts
from independently held-out natural images. If
$\x_0\sim\mathcal N(0,\Sigma)$ independently of the fit and the step, then,
conditional on the fit,
\[
 \widetilde{\x}_0=\x_0+m\bhat
 \sim\mathcal N(m\bhat,\Sigma).
\]
Unless $m\bhat=0$, this is not the original natural distribution. Unconditionally
the shifted seed is also coupled to training through $\bhat$ (and through $m$
if the schedule is fit-dependent). It cannot be treated as a fresh independent
natural seed in an RMT substitution. For a general natural-image distribution,
translation likewise need not preserve the distribution or its support.
Thus centering $\alpha$ is always possible algebraically, but imposing both
centered steps and independent seeds from the original natural distribution
restricts the protocol. The mean-bias term in Eq.~\eqref{eq:single} is retained
in the shifted baseline; it has not disappeared.

For seed-dependent schedules one may instead center conditionally:
$\widetilde\alpha=\alpha-\E[\alpha\mid\x_0,\bhat]$ and
$\widetilde{\x}_0=\x_0+\E[\alpha\mid\x_0,\bhat]\bhat$.
Then $\E[\widetilde\alpha\mid\x_0,\bhat]=0$ and hence the residual step is
uncorrelated with the transformed seed conditional on the fit, but generally
not independent of it. This pushforward of the seed distribution can change
both its mean and covariance. The master moment formula below still applies.
For natural-image experiments we therefore keep the original natural seed
and explicitly retain the schedule mean; for algebraic analysis either
parameterization is valid.

\paragraph{Student-response normalization.}
Let $u=D\alpha$, $\E u=\mu$, and $\Var u=\tau^2$.
Choosing $\tau^2=S=(\bstar)^\top\Sigma\bstar$ reproduces the old range convention.
Then
\begin{align}
 E&=\tau^2(1-g)^2+[\delta_0+(1-g)\mu]^2,\label{eq:single-normalized}\\
 R^2&=1-\frac{(1-g)^2}{g^2}
       -\frac{[\delta_0+(1-g)\mu]^2}{g^2\tau^2}.
\end{align}
Variance matching alone no longer fixes MSE when $\mu\ne0$.
For increment scoring, $\Delta P=u$, $\Delta Y=g u$, so
\[
 E_\Delta=(\tau^2+\mu^2)(1-g)^2,\qquad
 R^2_\Delta=1-\frac{(\tau^2+\mu^2)(1-g)^2}{g^2\tau^2}.
\]
Subtracting each path's own mean instead gives centered loss
$\tau^2(1-g)^2$ and the old $R^2=1-(1-1/g)^2$.
Absolute-response, increment, and centered-response scoring answer different
questions. The previous formulas for arbitrary $\x_0$ require the latter
convention, or suitable zero-baseline and zero-mean assumptions.

\section{Master formula for multiple seeds and dependent schedules}
Let $(\x_0,\alpha)$ have any joint evaluation distribution, with finite
second moments conditional on the fitted model. Define
\[
 \mu_0=\E\x_0,\quad \Sigma_0=\Cov(\x_0),\quad
 m=\E\alpha,\quad v=\Var\alpha,\quad
 \bm c_{0\alpha}=\Cov(\x_0,\alpha).
\]
Here $\bm c_{0\alpha}$ is the seed--step covariance. The generated input mean and covariance are exactly
\begin{equation}
 \mu_p=\mu_0+m\bhat,\qquad
 \Sigma_p=\Sigma_0+v\bhat\bhat^\top+\bm c_{0\alpha}\bhat^\top+\bhat\bm c_{0\alpha}^\top.
 \label{eq:input-covariance}
\end{equation}
Define $V_P=\bhat^\top \Sigma_p\bhat$, $V_Y=(\bstar)^\top \Sigma_p\bstar$, and
$C_{PY}=\bhat^\top \Sigma_p\bstar$. Then
\begin{equation}\boxed{
 E=\bm e^\top \Sigma_p\bm e+(\delta c+\bm e^\top\mu_p)^2,\quad
 R^2=1-E/V_Y,\quad \mathrm{slope}_p=C_{PY}/V_P,\quad
 r_p=C_{PY}/\sqrt{V_PV_Y}.}
 \label{eq:master}
\end{equation}
Explicitly,
\begin{align}
 V_P&=\bhat^\top \Sigma_0\bhat+vD^2+2D\bhat^\top\bm c_{0\alpha},\\
 V_Y&=(\bstar)^\top \Sigma_0\bstar+vN^2+2N(\bstar)^\top\bm c_{0\alpha},\\
 C_{PY}&=\bhat^\top \Sigma_0\bstar+vDN+N\bhat^\top\bm c_{0\alpha}+D(\bstar)^\top\bm c_{0\alpha},\\
 E&=\bm e^\top \Sigma_0\bm e+vL^2+2L\bm e^\top\bm c_{0\alpha}
       +(\delta c+\bm e^\top\mu_0+mL)^2.
\end{align}
This includes adaptive schedules and seed selection; no Gaussian evaluation
assumption is needed. However, $\Sigma_0$ or $\bm c_{0\alpha}$ depending on the training fit
cannot automatically be treated as deterministic weights in an RMT limit.

\paragraph{Finite seed collection.}
For seeds $\x_j$ with weights $\pi_j$ summing to one, and conditional
step means $m_j$ and variances $v_j$, use
\begin{align}
 \mu_0&=\sum_j\pi_j\x_j,&
 \Sigma_0&=\sum_j\pi_j(\x_j-\mu_0)(\x_j-\mu_0)^\top,\\
 m&=\sum_j\pi_jm_j,&
 v&=\sum_j\pi_j[v_j+(m_j-m)^2],\\
 \bm c_{0\alpha}&=\sum_j\pi_j(\x_j-\mu_0)(m_j-m).
\end{align}
With empirical step moments and observation-count weights, these formulas
are exact for the finite pooled data (use variance divisor equal to the
number of observations). Equal seed weighting differs from equal image
weighting if paths have different lengths. Pooled $R^2$, slope, and
correlation are generally not averages of per-seed metrics. All seeds share
one fitted weight: adding seeds reduces evaluation sampling error, not
training-fit uncertainty.

\section{Independent seeds: the prediction--control mixture}
Suppose $u=D\alpha$ is independent of the seed conditional on the fit, with
mean $\mu$ and variance $\tau^2$. Introduce the seed moments
\[
 D_0=\bhat^\top \Sigma_0\bhat,\quad N_0=\bhat^\top \Sigma_0\bstar,\quad
 S_0=(\bstar)^\top \Sigma_0\bstar,\quad
 \delta_{\rm seed}=\delta c+\bm e^\top\mu_0.
\]
Equation~\eqref{eq:master} reduces to
\begin{align}
 E&=D_0-2N_0+S_0+\tau^2(1-g)^2
           +[\delta_{\rm seed}+(1-g)\mu]^2,\label{eq:ind-mse}\\
 R^2&=1-\frac{E}{S_0+g^2\tau^2},\\
 \mathrm{slope}_p&=\frac{N_0+g\tau^2}{D_0+\tau^2},\qquad
 r_p=\frac{N_0+g\tau^2}
 {\sqrt{(D_0+\tau^2)(S_0+g^2\tau^2)}}.\label{eq:ind-corr}
\end{align}
When $D_0>0$, the pooled slope is the variance-weighted average
\[
 \mathrm{slope}_p=\frac{D_0}{D_0+\tau^2}\frac{N_0}{D_0}
       +\frac{\tau^2}{D_0+\tau^2}g.
\]
The normalization $\tau^2=S$ matches the \emph{within-seed student variance};
the pooled student variance is $D_0+S$, not $S$.
If instead total student variance must equal $S$, one needs
$\tau^2=S-D_0\geq0$ under this independent protocol.

For centered natural seeds, $\Sigma_0=\Sigma$, $\mu_0=0$, zero intercepts,
and centered increments, let $E_{\rm gen}=D_0-2N_0+S$. Then
\begin{equation}\boxed{
 E=E_{\rm gen}+\tau^2(1-g)^2,\qquad
 R^2_{\rm pooled}=1-\frac{E_{\rm gen}+\tau^2(1-g)^2}{S+g^2\tau^2}.}
\end{equation}
At $\tau^2=S$, this is
$R^2_{\rm pooled}=1-[E_{\rm gen}/S+(1-g)^2]/(1+g^2)$.
Thus multiple natural seeds bring ordinary generalization error back into
accentuation evaluation. At $\tau^2=0$, all metrics are seed-prediction
metrics. For fixed $g\ne0$ and $\mu=0$, as $\tau^2$ dominates seed variance,
slope tends to $g$, Pearson correlation tends to $\operatorname{sign}(g)$,
and $R^2$ tends to the single-path value. At $g=0$, seed teacher variance
can keep pooled $R^2$ defined even though per-path $R^2$ is undefined.

\section{A common alternative: absolute student target levels}
Suppose the experiment selects a target response $q$, independently of seeds,
and moves until $P(\x)=q$. Then
\[
 \alpha=(q-p_0)/D,\qquad P=q,\qquad
 Y=g q+a_0,\qquad a_0=y_0-g p_0.
\]
This makes the step seed-dependent, even for a centered target distribution.
Let $\E q=\nu$, $\Var q=t^2>0$, $\E a_0=\bar a$,
and $\Var a_0=v_a$. Direct calculation gives
\begin{align}
 E&=(1-g)^2t^2+v_a+[(1-g)\nu-\bar a]^2,\\
 R^2&=1-\frac{E}{g^2t^2+v_a},\qquad \mathrm{slope}_p=g,\\
 r_p&=\frac{g t}{\sqrt{g^2t^2+v_a}},\qquad
 v_a=(\bstar-g\bhat)^\top \Sigma_0(\bstar-g\bhat).
\end{align}
Here $\bm c_{0\alpha}=-\Sigma_0\bhat/D$. Unlike independent increments, the pooled slope
remains $g$, while Pearson correlation measures the variation of path
intercepts as well. Consequently the experimental schedule must be specified
before quoting a multi-seed metric formula.

\section{A specific protocol: natural seeds and natural response matching}
\label{sec:natural-matching}
Take independent held-out natural seeds with mean $\bm m_{\rm nat}$ and
covariance $\Sigma$, and condition throughout on one fitted model. Define
\[
 \nu_P=\widehat c+\bhat^\top\bm m_{\rm nat},\qquad
 \nu_Y=c_\star+(\bstar)^\top\bm m_{\rm nat},\qquad
 \delta=\nu_P-\nu_Y,
\]
\[
 D_{gen}=\bhat^\top\Sigma\bhat,\quad N_{gen}=\bhat^\top\Sigma\bstar,\quad
 S=(\bstar)^\top\Sigma\bstar>0,\quad g=N/D.
\]
The reference distribution is the \emph{noiseless natural teacher response},
with mean $\nu_Y$ and variance $S$. We first require the generated
\emph{student absolute response} to match these moments. This is the
nonzero-seed analogue of the original student-range normalization; it does
not assume that the generated teacher response is already calibrated.
Matching two moments implies full distributional equality only in a Gaussian
setting, or when we explicitly sample target responses from the reference law.

\subsection{Match the natural response range separately from each seed}
For every seed, draw an absolute target $q$ independently from the natural
teacher response distribution and set
\begin{equation}\boxed{
 \alpha=\frac{q-P(\x_0)}{D}.}
 \label{eq:matching-target}
\end{equation}
Then $P(\x_0+\alpha\bhat)=q$ exactly, so the entire student response
distribution matches the reference, conditionally on \emph{each} seed.
In particular,
\begin{equation}\boxed{
 \E[\alpha\mid\x_0,\bhat]=\frac{\nu_Y-p_0}{D},\qquad
 \Var(\alpha\mid\x_0,\bhat)=\frac{S}{D^2}.}
 \label{eq:matching-conditional}
\end{equation}
For Gaussian natural inputs these conditional steps are Gaussian. For a
non-Gaussian response reference, use its actual quantiles or samples for $q$;
mean and variance alone do not specify its law.

A seed with a high predicted baseline needs a negative mean step; one with a
low predicted baseline needs a positive mean step. Even if the global natural
response is centered, the step is generally not centered for an individual
seed. Averaging over natural seeds gives, by total expectation and variance,
\begin{equation}\boxed{
 \E\alpha=-\frac{\delta}{D},\qquad
 \Var\alpha=\frac{S+D_{gen}}{D^2},\qquad
 \Cov(\x_0,\alpha)=-\frac{\Sigma\bhat}{D}.}
 \label{eq:matching-marginal}
\end{equation}
The marginal step variance is larger than the per-seed variance: it includes
the variation of the baseline-canceling conditional means. For centered
zero-intercept models, $\E\alpha=0$ marginally but
$\E[\alpha\mid\x_0]=-p_0/D$ is usually nonzero.
Ignoring this seed--step anticorrelation gives an incorrect pooled variance.
Indeed,
\[
 \Var P_{\rm acc}=D_{gen}+D^2\Var\alpha
                  +2D\Cov(p_0,\alpha)=D_{gen}+(S+D_{gen})-2D_{gen}=S.
\]

Conditionally centering the step produces the shifted seed
\[
 \widetilde{\x}_0=\x_0+\frac{\nu_Y-p_0}{D}\bhat,
 \qquad P(\widetilde{\x}_0)=\nu_Y,
\]
with mean $\bm m_{\rm nat}-\delta\bhat/D$ and covariance
$J\Sigma J^\top$, where $J=I-\bhat\bhat^\top/D$.
Thus the shifted seeds lie on a common student-response hyperplane and are
not independent samples from the original natural distribution. The original
seeds are natural; the recentered mathematical seeds need not be.

\subsection{Metrics and the remaining RMT quantity}
Let $a_0=y_0-g p_0$ and
\begin{equation}\boxed{
 A_{\rm seed}=\Var(a_0)
  =(\bstar-g\bhat)^\top\Sigma(\bstar-g\bhat)
  =S-2g N_{gen}+g^2D_{gen}.}
 \label{eq:matching-residual}
\end{equation}
The teacher response on the generated images is $Y_{\rm acc}=g q+a_0$.
Since $q$ and $a_0$ are independent,
\begin{align}
 \E Y_{\rm acc}&=\nu_Y-g\delta,&
 \Var Y_{\rm acc}&=g^2S+A_{\rm seed},\\
 E&=(1-g)^2S+A_{\rm seed}+g^2\delta^2,\label{eq:matching-mse}\\
 R^2&=1-\frac{(1-g)^2S+A_{\rm seed}+g^2\delta^2}
                    {g^2S+A_{\rm seed}},\\
 \mathrm{slope}_p&=g,&
 r_p&=\frac{g\sqrt S}{\sqrt{g^2S+A_{\rm seed}}}.
\end{align}
The denominators must be nonzero. In the centered case, the compact dimensionless
form is
\[
 \frac{E}{S}=(1-g)^2+\frac{A_{\rm seed}}S,\qquad
 R^2=1-\frac{(1-g)^2+A_{\rm seed}/S}{g^2+A_{\rm seed}/S},\qquad
 r_p=\frac{g}{\sqrt{g^2+A_{\rm seed}/S}}.
\]
Thus natural-response matching leaves two distinct sources of error: path gain
miscalibration and variation of the teacher response across equal-student-response
seed locations. Even though student responses have the correct marginal law,
the teacher law generally does not match it.
The pooled slope is still $g$, but Pearson correlation is no longer forced to
be $\pm1$. The additional RMT summary is $A_{\rm seed}/S$, constructed from
the existing natural-prediction moments $N_{gen},D_{gen}$ and the control ratio $g$.
A leading concentrated-fit DE uses
\[
 A_{\rm seed}^{\rm DE}=S-2g_{\rm DE}N_{gen,DE}
                         +g_{\rm DE}^2D_{gen,DE}.
\]
Its fit average instead requires $\E[g N_{gen}]$ and $\E[g^2D_{gen}]$, so joint fluctuations
matter when these quantities do not concentrate. As a check, if
$\bhat=k\bstar$ with $k\ne0$ and zero intercepts, $g=1/k$ and
$A_{\rm seed}=0$: multi-seed correlation is still $\operatorname{sign}(k)$,
while gain miscalibration remains. At $\bhat=\bstar$, all metrics are perfect even
though the protocol can move images substantially.

\subsection{Pooled matching alone does not identify a unique schedule}
Suppose we only require the \emph{pooled} student responses to have mean
$\nu_Y$ and variance $S$. Put $c_{p\alpha}=\Cov(p_0,\alpha)$. Then
\begin{equation}
 \E\alpha=-\delta/D,\qquad
 D^2\Var\alpha+2D c_{p\alpha}=S-D_{gen}.
 \label{eq:pooled-matching-constraint}
\end{equation}
Without specifying dependence on the seed, the variance is not uniquely
identified. Under \emph{independent steps}, $c_{p\alpha}=0$, hence
\[
 \Var\alpha=(S-D_{gen})/D^2,
\]
which is feasible only for $D_{gen}\leq S$. In the centered case this adds only the
missing student variance; if $D_{gen}=S$, it forces no motion. If $D_{gen}>S$, independent
steps cannot reduce response variance. Independent Gaussian steps give exact
Gaussian response matching when natural seeds are Gaussian, but only pooled
matching, not matching separately from each seed.

More generally, couple a target $q$ with the seed while keeping its desired
marginal distribution, and let $c=\Cov(q,p_0)$. Equation~\eqref{eq:matching-target}
gives
\[
 \Var\alpha=\frac{S+D_{gen}-2c}{D^2},\qquad |c|\leq\sqrt{S D_{gen}}.
\]
The independent-target protocol above chooses $c=0$. For Gaussian seeds with
$D_{gen}>0$, the deterministic choice
$q=\nu_Y+\sqrt{S/D_{gen}}(p_0-\nu_P)$ gives the same Gaussian pooled reference law
with $c=\sqrt{S D_{gen}}$ and the smallest possible step variance
$(\sqrt S-\sqrt{D_{gen}})^2/D^2$. It yields one point per seed rather than a full
response range from every seed. The metrics in the previous subsection assume
independent targets and must not be reused for this correlated-target protocol.

\subsection{Which response is being matched?}
Matching generated student responses to \emph{natural student} responses
would use $(\nu_P,D_{gen})$ instead of $(\nu_Y,S)$. Independent-step pooled matching
then forces $\alpha=0$ almost surely, whereas independent absolute targets
still allow motion, with conditional variance $D_{gen}/D^2$ and marginal variance
$2D_{gen}/D^2$.

If the requirement instead concerns the \emph{true teacher} response after
accentuation, its moments obey
\[
 \E Y_{\rm acc}=\nu_Y+N\E\alpha,\qquad
 \Var Y_{\rm acc}=S+N^2\Var\alpha+2N\Cov(y_0,\alpha).
\]
For $N\ne0$, matching the natural teacher moments with seed-independent steps
forces $\E\alpha=0$ and $\Var\alpha=0$. Nontrivial matching therefore needs
seed-dependent steps, for example the oracle schedule
$\alpha=(q-y_0)/N$ with an independent natural teacher target $q$.
This is not the student-controlled protocol and requires teacher information.
At $N=0$, every step leaves teacher responses unchanged, so teacher-distribution
matching imposes no constraint on motion.

Finally, exact moment matching can conflict with purely one-sided
accentuation ($\alpha\geq0$). In the centered, equal-mean case it requires
$\E\alpha=0$, which under that constraint forces $\alpha=0$ almost surely.
A nontrivial matched-range experiment generally includes both accentuation
and de-accentuation, or deliberately uses a shifted reference distribution.

\section{The key RMT quantities}
\subsection{Which scalars are sufficient?}
For independent increments, the conditional metrics depend on
\begin{equation}\boxed{
 (N,D,\ N_0,D_0,\ \ell_0=\bhat^\top\mu_0),}
 \label{eq:core}
\end{equation}
together with deterministic teacher moments, intercepts, and schedule
moments. A single fixed seed has $\Sigma_0=0$, so the new object is
$\ell_0=\bhat^\top\x_0$, or equivalently $\delta_0$.
For centered natural seeds, $N_0,D_0$ are exactly the existing
$N_{\rm gen},D_{\rm gen}$; no new leading-order scalar theory is needed.
For general seeds, the new geometry is $\Sigma_0$, not necessarily $\Sigma$ or $I$.
For dependent schedules add $\bhat^\top\bm c_{0\alpha}$ and $(\bstar)^\top\bm c_{0\alpha}$;
when the schedule depends on the fit, express these through the actual
protocol before applying a DE.

\subsection{Ridge resolvents and leading deterministic equivalents}
For direct pixel ridge let
\[
 \hat\Sigma=X^\top X/n,\quad Q=(\hat\Sigma+\lambda I)^{-1},\quad
 \bhat=(I-\lambda Q)\bstar+\frac{\sigma}{n}QX^\top\eta.
\]
Conditional on $X$, write $\bhat=\bm m_X+\bm\xi$ with
$\bm m_X=(I-\lambda Q)\bstar$ and
$\Omega_X=(\sigma^2/n)Q\hat\Sigma Q$. For deterministic $\Sigma_0$ and
$\bm a$, the exact conditional identities are
\[
 \E_\eta[\bm a^\top\bhat\mid X]=\bm a^\top\bm m_X,\qquad
 \E_\eta[\bhat^\top \Sigma_0\bhat\mid X]
 =\bm m_X^\top \Sigma_0\bm m_X+\Tr(\Sigma_0\Omega_X).
\]
Thus the required random-matrix objects are one-resolvent bilinear forms
$\bm a^\top Q\bstar$, weighted two-resolvent forms
$(\bstar)^\top Q\Sigma_0 Q\bstar$, and traces $\Tr(\Sigma_0 Q\hat\Sigma Q)$.
Replacing $Q$ independently in each factor of $Q\Sigma_0 Q$ loses design variance.

Using the same leading aggregate DE as the existing feature note, set
\begin{align}
 \kappa-\lambda&=\frac{\kappa}{n}\Tr[\Sigma(\Sigma+\kappa I)^{-1}],\\
 H&=(\Sigma+\kappa I)^{-1},\qquad
 \tilde\beta_\kappa=\Sigma H\bstar,\\
 v_0&=\frac{\kappa^2(\bstar)^\top H\Sigma H\bstar+\sigma^2}
 {n-\Tr(\Sigma^2H^2)},\qquad V_{\beta,DE}=v_0H\Sigma H.
\end{align}
In the main manuscript's spectral notation,
\[
 B_{a,b}(\kappa)=\sum_k\frac{s_k^a b_k^2}{(s_k+\kappa)^b},\qquad
 \mathrm{df}_{a,b}(\kappa)=\sum_k\frac{s_k^a}{(s_k+\kappa)^b},\qquad
 \mathrm{df}_2=\mathrm{df}_{2,2}.
\]
Hence $v_0=(\kappa^2B_{1,2}+\sigma^2)/(n-\mathrm{df}_2)
=(E_{gen,DE}+\sigma^2)/n$. The manuscript's training-fit means
$\mu_N=\E_{X,\eta}N$ and $\mu_D=\E_{X,\eta}D$ have leading equivalents
$N_{\rm DE}=B_{1,1}$ and
$D_{\rm DE}=B_{2,2}+v_0\mathrm{df}_{1,2}$, respectively;
these DEs are not exact finite-sample means.

Then the aggregate substitutions are
\begin{align}
 N_{\rm DE}&=(\bstar)^\top\tilde\beta_\kappa,&
 D_{\rm DE}&=\|\tilde\beta_\kappa\|^2+\Tr V_{\beta,DE},\\
 N_{0,DE}&=\tilde\beta_\kappa^\top \Sigma_0\bstar,&
 D_{0,DE}&=\tilde\beta_\kappa^\top \Sigma_0\tilde\beta_\kappa+\Tr(\Sigma_0 V_{\beta,DE}),\\
 \ell_{0,\rm DE}&=\mu_0^\top\tilde\beta_\kappa,&
 g_{\rm DE}&=N_{\rm DE}/D_{\rm DE}.
\end{align}
These are leading aggregate equivalents under the same fixed-geometry,
boundedness, and concentration assumptions as the source note. $V_{\beta,DE}$ is an
aggregate second-moment prescription, not an exact finite-dimensional
covariance or a full joint CLT. In particular, for a fixed finite number of
seeds, baseline projections need not concentrate on the scale relevant to
their squared errors. When they do not, retain the exact conditional
Gaussian moments above or a justified joint fluctuation approximation.

In a common eigenbasis of $\Sigma_0$ and $\Sigma$, with eigenvalues $s_{0,k}$
and $s_k$, and teacher coordinates $b_k$, the new quadratic form reads
\[
 D_{0,DE}=
 \sum_k s_{0,k}\left[
 \left(\frac{s_kb_k}{s_k+\kappa}\right)^2
 +v_0\frac{s_k}{(s_k+\kappa)^2}\right].
\]
This shows exactly how seed variability weights the same spectral quantities.
Noncommuting $\Sigma_0$ requires the matrix form, not just its diagonal spectrum.

\subsection{Linear features}
Use the existing notation $\Sigma_F=F\Sigma F^\top$, $G=FF^\top$,
$\bm q=F\bstar$, and population feature coefficient $w^\star$.
With $m_{w,DE}$ and $V_{w,DE}$ from the feature-space DE, replace
\[
 \tilde\beta_\kappa=F^\top m_{w,DE},\qquad V_{\beta,DE}=F^\top V_{w,DE} F.
\]
Control uses $G$; seed prediction adds $\Sigma_{0,F}=F\Sigma_0 F^\top$ and
$\bm q_0=F\Sigma_0\bstar$. In particular,
\[
 N_{0,DE}=m_{w,DE}^\top\bm q_0,\qquad
 D_{0,DE}=m_{w,DE}^\top \Sigma_{0,F}m_{w,DE}+\Tr(\Sigma_{0,F}V_{w,DE}).
\]
Thus training covariance, control geometry, and seed covariance are three
potentially different matrices. The teacher residual outside the feature
span is already included by retaining the full teacher $S_0$ and $\bm q_0$.

\subsection{Averaging over fitted models needs joint fluctuations}
Even the one-seed normalized MSE is not determined by $\E g$ alone.
For fixed $\mu,\tau^2$, exactly
\begin{align}
 \E_{\rm fit}E
 &=\tau^2\{(1-\bar g)^2+\Var_{\rm fit}g\}
   +[\bar\delta_0+\mu(1-\bar g)]^2\\
 &\quad+\Var_{\rm fit}\delta_0+\mu^2\Var_{\rm fit}g
       -2\mu\Cov_{\rm fit}(\delta_0,g).
\end{align}
Thus the new fluctuation object is the \emph{joint} law of baseline error and
control gain, not only $\Var(g)$. Writing
$\bm j_g=(\bstar-2g\bhat)/D$ gives the local differential $dg=\bm j_g^\top d\bhat$.
At fixed $X$, a first-order Gaussian delta calculation around
$\bm m_X$ (where its norm is nonzero) gives
\[
 \Var_\eta(\delta_0\mid X)=\x_0^\top\Omega_X\x_0,
 \qquad
 \Cov_\eta(\delta_0,g\mid X)\simeq
 \x_0^\top\Omega_X\bm j_g\big|_{\bhat=\bm m_X}.
\]
This assumes fixed intercepts and small conditional fluctuations. Full
unconditional covariance also contains variation of conditional means over
$X$, as in the existing variance note. Pooled ratio metrics require the joint
law of Eq.~\eqref{eq:core}; CV additionally randomizes $\lambda$.
Neither ratios of means nor a Gaussian fitted-weight approximation are exact
expectations of nonlinear metrics. Near $N=0$ or a small true-response
variance, report distributions/quantiles and assess denominator concentration.

\section{Practical reporting and limits}
For each experiment specify: absolute responses versus increments versus
within-path centering; independent increments versus absolute targets;
step mean and variance; seed/observation weights; and pooled versus per-seed
metrics. Report MSE and uncalibrated $R^2$ alongside slope and Pearson $r_p$.
A high Pearson $r_p$ on a single linear path establishes monotonicity only.
Pooling seed baselines can also yield high correlation even when the control
gain is poor; within-seed gains make that failure visible.

If independent evaluation noise of variance $\sigma_{\rm test}^2$ is added
to each teacher response, add it to both MSE and $V_Y$, leaving $C_{PY}$ and
$V_P$ unchanged. Slope is unchanged, correlation is attenuated, and the
noiseless $\pm1$ single-path result no longer holds. Subtracting a noisy seed
response creates shared noise across increments and needs its covariance
explicitly included.

All exact results assume a common, constant gradient $\bhat$. Nonlinear
networks, clipping, or image constraints generally make gains seed- and
step-dependent. A local linearization can apply per seed, but a pooled
analysis then needs the distribution of those local gains and their
covariance with baselines; the single ratio $N/D$ is no longer sufficient.
\end{document}
